\documentclass{article}
\usepackage[utf8]{inputenc}
\usepackage{amsmath}
\usepackage{geometry}
\geometry{margin=2cm}
\begin{document}

\section*{Análise da Questão 98 — ENEM 2024 — Ciências da Natureza}

\subsection*{Interpretação do enunciado e estratégia de resolução}

A questão apresenta uma situação em que placas indicadoras de saída de emergência possuem brilho no escuro devido à presença de sulfeto de zinco (ZnS) dopado com prata ou cobre. O fenômeno físico associado é a fosforescência, que resulta em emissão luminosa prolongada após exposição à luz ambiente.

\paragraph{Estratégia}
Para resolver a questão, deve-se:
\begin{itemize}
  \item Compreender o fenômeno descrito (fosforescência).
  \item Relacionar o processo às transições eletrônicas internas do material dopado.
  \item Avaliar as alternativas descartando conceitos incorretos (colisões, coloração, reações nucleares e reflexão).
  \item Selecionar a alternativa correta que explique o fenômeno como emissão luminosa por transições eletrônicas.
\end{itemize}

\subsection*{Conteúdo teórico e revisão}

\paragraph{Fosforescência}
Fenômeno físico-químico onde certos materiais absorvem energia luminosa e a liberam lentamente em forma de luz, mesmo após cessar a fonte de excitação. Difere da fluorescência pela duração da emissão.

\paragraph{Transições eletrônicas}
Mudanças nos níveis de energia de elétrons em átomos ou moleculas que resultam em absorção ou emissão de fótons.

\paragraph{Dopagem de materiais}
Incorporação de pequenas quantidades de átomos de outros elementos que alteram propriedades ópticas e elétricas, no caso favorecendo fosforescência.

\subsection*{Resolução e "Pulo do Gato"}

Ao analisar o fenômeno do brilho, nota-se que:
\begin{itemize}
  \item Emissão luminosa ocorre após a fonte de luz ambiente ser removida.
  \item O processo está relacionado a elétrons que ficam em estados excitados temporariamente.
  \item Outras alternativas que envolvem colisões, cor dos átomos, reações nucleares ou reflexão são incompatíveis com o fenômeno apresentado.
\end{itemize}

\paragraph{Pulo do Gato}
\textit{Fosforescência é luz emitida após absorver energia, por transições eletrônicas que retardam a emissão; não é reflexão nem processo nuclear.}

\subsection*{Armadilhas e justificativas das alternativas erradas}

\begin{itemize}
  \item \textbf{Alternativa A:} Brilho por colisões entre átomos — atrai por confundir energia transferida com emissão luminosa, porém está incorreta para a fosforescência.
  \item \textbf{Alternativa B:} Brilho pela cor dos átomos — confunde cor aparente com emissão luminosa real.
  \item \textbf{Alternativa D:} Brilho gerado por reações nucleares — incorreto pelo tipo de fenômeno e escala energética.
  \item \textbf{Alternativa E:} Brilho por reflexão — descartado pois a luz emitida ocorre após o término da luz ambiente.
\end{itemize}

\subsection*{Código da Questão}

\texttt{CN\_FISICA\_FENOMENOS\_001}


\end{document}